% =====================================================
% Two-channel decomposition: how a tax moves a firm's
% optimal ad spend. The margin channel and the demand
% rotation channel; pass-through arbitrates between them.
% =====================================================

\documentclass[11pt]{article}
\usepackage[margin=1in]{geometry}
\usepackage{amsmath,amssymb}
\usepackage{booktabs}
\usepackage{parskip}
\usepackage{hyperref}
\begin{document}

\section*{Two Channels From Tax To Ads}

\paragraph{Notation.} Demand for the firm's product is
$D(p, A)$, a function of retail price $p$ and ad spend
$A$. The own-price elasticity of demand is
\[
    \eta_p \;\equiv\; -\frac{\partial \ln D}{\partial \ln p},
\]
unitless and positive: a 1\% price increase cuts demand
by $\eta_p$ percent. $\eta_p > 1$ is elastic, $\eta_p < 1$
inelastic. Marginal cost to the firm is $c$, and
\[
    m \;\equiv\; \frac{p - c}{p}
\]
is the firm's per-unit profit as a fraction of the
consumer-facing price. For CPG manufacturers, this is
roughly the wholesale-gross-profit-per-retail-dollar.

\paragraph{Setup.} A specific tax $\tau$ is levied on
the firm per unit sold. Let $\rho \in [0,1]$ be the
share passed through to consumers, so retail price
rises by $\rho\tau$ and the firm's per-unit margin
falls by $(1-\rho)\tau$. The firm chooses $A$ to
maximize
\[
    \pi(A) \;=\;
    \big(p^{\text{firm}} - c\big)\,
    D\!\big(p^{\text{retail}},\, A\big)
    \;-\; A,
\]
with $p^{\text{firm}} = p^{\text{retail}} - \tau$. The
ad FOC is $(p^{\text{firm}}-c)\, D_A = 1$.

\paragraph{Assumptions used to derive the
decomposition below.}
\begin{enumerate}
    \item \textbf{Firm chooses $A$ only.} Retail price
    is set by the retailer or market; pass-through
    $\rho$ is exogenous. (Rules out the joint
    Dorfman--Steiner monopolist case, where $\rho$ is
    optimized and the margin channel is absorbed into
    the pricing decision.)
    \item \textbf{Separable, constant-elasticity
    demand:} $D(p, A) = D_0(p)\, A^{\eta_A}$ with
    $D_0(p) = p^{-\eta_p}$. Implies
    $\partial \eta_p / \partial \ln A = 0$
    (separability), which we relax in a footnote
    below.
    \item \textbf{Small $\eta_A$.} The marketing
    literature reports ad-spend elasticities of demand
    in $[0.05, 0.15]$, so $1/(1-\eta_A) \approx 1$ to
    within $\sim 5\text{--}15\%$.
\end{enumerate}

\paragraph{The two channels.} Differentiating the FOC
around the pre-tax optimum under these assumptions, and
keeping leading-order terms in $\tau$:
\begin{equation}\label{eq:two-channel}
    \boxed{\;
    \Delta \ln A^* \;\approx\;
    \underbrace{-\,(1-\rho)\,\frac{\tau}{p-c}}_{\substack{\text{margin}\\\text{channel}}}
    \;-\;
    \underbrace{\rho\,\eta_p\,\frac{\tau}{p}}_{\substack{\text{demand rotation}\\\text{channel}}}\,.
    \;}
\end{equation}
The margin channel: every unit of tax the firm absorbs
cuts marginal profit per ad-acquired customer
one-for-one, so ads fall in proportion to the
tax-as-share-of-margin $\tau / (p-c)$.

The demand rotation channel: the passed-through tax
raises consumer price, demand falls by $\eta_p \rho
\tau / p$, the marginal product of advertising falls
proportionally (since $D_A = \eta_A D / A$ under
separable demand), and the firm cuts $A$ to restore the
FOC.

Both channels are unambiguously negative: under
separable demand at the assumed parameters, a tax cuts
ads through both routes.

\paragraph{Where the literature's $(1-\eta_p)$ form
comes from, and why we are not using it.} Standard
textbook Dorfman--Steiner derives $A^* = (\eta_A /
\eta_p)\, pQ$ for a monopolist that jointly optimizes
$p$ and $A$, implying $\Delta \ln A^* = (1-\eta_p)\Delta
\ln p$. That result requires the firm to be repricing
optimally, which forces $\rho$ to equal the
profit-maximizing pass-through $\eta_p / (\eta_p - 1)$
under constant-elasticity demand (over 100\% for
inelastic markups) and absorbs the margin response
into the pricing decision. Empirical pass-through for
CPG excise taxes runs $\rho \in [0.7, 1.0]$, not the
D--S monopolist optimum, so we use the $A$-only
response above instead --- which keeps the margin
channel additive at any exogenous $\rho$.

\paragraph{Non-separable correction.} Relaxing
$\partial \eta_p / \partial \ln A = 0$ adds a term
$-\rho\,(\tau/p)\,\eta_A^{-1}\, \partial \eta_p /
\partial \ln A$ to the demand rotation channel. The
amplification factor $1/\eta_A \approx 20$ for $\eta_A
= 0.05$ means that even a small elasticity-shift would
move $\Delta \ln A^*$ substantially. This is the
``persuasion-rotates-demand'' object we cannot identify
from the data without a calibration target; we note it
but treat the separable case as the working model.

\paragraph{Pass-through arbitrates the two channels.}
\eqref{eq:two-channel} has a clean interpretation as a
function of $\rho$:
\begin{itemize}
    \item $\rho = 0$ (no pass-through): only the margin
    channel operates. Response is $-\tau/(p-c) =
    -(\tau/p)/m$, scaling with tax-as-share-of-margin.
    First order in $\tau$, amplified by the inverse
    margin $1/m$.

    \item $\rho = 1$ (full pass-through): only the
    demand rotation channel operates. Response is
    $-\eta_p \tau/p$. First order in $\tau$, scaled by
    the price elasticity of demand.

    \item Intermediate $\rho$: both channels add
    (both negative under separable demand).
\end{itemize}

\paragraph{Ratio of the two channels at their
respective extremes.} Comparing full pass-through to no
pass-through, holding the tax fixed:
\[
    \frac{|\,\text{demand rotation channel at } \rho=1\,|}
         {|\,\text{margin channel at } \rho=0\,|}
    \;=\;
    \eta_p \cdot \frac{p-c}{p}
    \;=\;
    \eta_p \cdot m.
\]
For the candidate categories ($m \in [0.4, 0.5]$,
$\eta_p \in [0.4, 0.6]$), this ratio falls in
$[0.16, 0.30]$ --- the demand rotation channel under
full pass-through is one sixth to one third the size of
the margin channel under zero pass-through.

\paragraph{Numbers for current candidates.} Predicted
$\hat R$ from \eqref{eq:two-channel}:

\begin{center}
\begin{tabular}{lcccccc}
\hline
Setting & $\rho$ & $\tau/p$ & $m$ & $\eta_p$ &
\multicolumn{2}{c}{$\hat R$ (signed)} \\
& & & & & margin & dem.\ rot. \\
\hline
SSB Philly                       & 1.00 & 0.13 & 0.50 & 0.60 & $\;\;0.000$ & $-0.078$ \\
Beer, state-level                & 0.85 & 0.01 & 0.50 & 0.50 & $-0.003$ & $-0.004$ \\
Cigarettes, US state             & 0.90 & 0.20 & 0.40 & 0.40 & $-0.050$ & $-0.072$ \\
\hline
\end{tabular}
\end{center}

\noindent Two remarks on this table:

\textbf{(i) SSB Philly predicts $-7.8\%$, not the
near-null we observe.} The literal A-only formula
produces an ad cut comfortably above the placebo
noise floor ($\sigma_{\hat R} \approx 0.05$). The most
likely reason we don't see it: ad spend is allocated
nationally and a city-level tax affects under a
percent of a brand's national demand, so the firm's
national FOC barely moves and DMA-level Spot TV in
Philly does not get reoptimized down by $\eta_p
(\tau/p)$. Adjustment frictions reinforce this: ad
budgets are set annually and reallocated coarsely. The
model as written assumes ads are reoptimized at the
same spatial scale as the tax, which is wrong for
sub-DMA city taxes.

\textbf{(ii) Cigarettes also above the noise floor.}
Under the same model, a 20\% cigarette tax with 90\%
pass-through predicts $\hat R \approx -12\%$. We
cannot test this --- cigarette broadcast ads are
banned in the US since 1971. Foreign or pre-1971
settings would.

\paragraph{Implications.}
\begin{enumerate}
    \item The literal model prediction for SSB Philly
    ($-7.8\%$) exceeds the observed null. The model
    assumes firm-level ad reoptimization at the spatial
    scale of the tax. For sub-DMA city taxes that
    affect under a percent of a national brand's
    demand, that assumption fails: ads are set
    nationally and don't reoptimize for local shocks.
    The formula is most credible at the spatial scale
    where the firm actually rebalances ads --- usually
    national for CPG.

    \item Margin and demand rotation channels both
    push ads down under separable demand. The
    structural object that could flip a sign is
    $\partial \eta_p / \partial \ln A$. It enters the
    demand rotation channel multiplied by $\rho\tau/p$
    and divided by $\eta_A$, so it is amplified by
    $\sim 20$ relative to the mechanical $-\eta_p$
    term --- which means even a small effect could
    matter if we had a calibration target.

    \item The contribution is the decomposition
    itself: DMA-level Spot TV ad response to a tax
    depends on $\rho$, $m$, $\eta_p$, and the spatial
    scale of firm reoptimization. None of these is
    innate to the category; they reflect institutional
    details that can be read off from industry data.
\end{enumerate}

\end{document}
